\clearpage\SourceLeaf{111}%
\pdfbookmark[1]{Über die Grenzen der geocentrischen Örter der Planeten}{grenzen1804}
\begin{center}
{\small ÜBER DIE GRENZEN DER GEOCENTRISCHEN ÖRTER DER PLANETEN.\par}
\medskip\hrule\medskip
{\small Monatliche Correspondenz zur Beförderung der Erd- und Himmels-Kunde,\par
herausgegeben vom Freiherrn \textsc{von Zach}. August 1804.\par}
\medskip\hrule
\end{center}\medskip
Von der Sonne aus gesehen erscheint die Bewegung jedes Planeten, in so
fern man auf die kleinen Störungen durch andere Himmelskörper nicht sieht,
stets in einem und demselben grössten Kreise am Fixsternhimmel. Eben so
würde sie auch von der Erde aus erscheinen, wenn die Ebene seiner Bahn mit
der Ecliptik zusammenfiele. Sind aber diese beiden Ebenen gegen einander geneigt,
so liegen alle mögliche geocentrische Örter des Planeten auf der Himmelskugel
nicht mehr, wie in jenem Falle, in einer nur nach einer Dimension ausgedehnten
Linie, sondern sie erfüllen einen Flächenraum, eine \textit{Zone}, die den
ganzen übrigen Himmel, wo der Planet von der Erde aus nie erscheinen kann,
in zwei Theile absondert, und füglich der \textit{Zodiacus} des Planeten heissen kann.
Auf diese Weise hat also jeder Planet im Grunde seinen eigenthümlichen Zodiacus,
dessen Grenzen (\textit{Limiten}) vollkommen scharf bestimmbar sind, in so fern
man seine und die Erd-Bahn als Kegelschnitte von unveränderlichen Elementen
ansieht. Die genaue Bestimmung dieser Grenzen ist an sich schon ein interessantes
analytisches Problem; aber die Anwendung desselben, besonders auf die
beiden neuen Planeten, deren Zodiacus eine beträchtliche Ausdehnung haben,
ist auch nicht ohne practische Wichtigkeit. Man weiss, dass zur Aufsuchung
und Beobachtung dieser merkwürdigen Himmelskörper sehr genaue und detaillirte
Sternkarten erfordert werden, und dass selbst die besten, welche wir bisher
\SourceLeaf{112}%
besitzen, dazu bei weiten noch nicht hinlänglich sind. Wenn man daher nicht
jedes Jahr von den Gegenden, die diese Planeten durchlaufen, Special-Karten
entwerfen will, so muss man nothwendig auf einen eigenen, den ganzen Raum,
worin sie sich zeigen können, begreifenden Atlas denken. Die genaue Bestimmung
der Grenzen dieses Raums wird daher um so wünschenswerther, da man
sich bei einem solchen Unternehmen, das an sich schon von bedeutendem Umfange
ist, gern alle zu diesem Zwecke unnöthige Mühe ersparen wird. Gewiss
wird allen Freunden der Astronomie die Nachricht sehr willkommen sein, dass
der geschickte Lilienthaler Astronom \textsc{Harding}, von dem wir bereits verschiedene
vortreffliche Specialkarten besitzen, schon angefangen hat, sich dieser
grössern Arbeit zu unterziehen, die sich nicht nur durch ein sehr reiches Detail,
sondern auch durch die sorgfältigste, durchgehends auf Autopsie gegründete Kritik
sehr vortheilhaft auszeichnen wird.

Wir legen durch den Mittelpunkt der Sonne drei auf einander senkrechte,
übrigens willkürliche Ebenen, und nennen die senkrechten Abstände des beobachteten
Planeten von denselben $x$, $y$, $z$; die Abstände der Erde hingegen
$x'$, $y'$, $z'$. Wir setzen ferner
\[
\begin{aligned}
x'-x&=\Delta\cos\alpha\cos\ObservedBeta\\
y'-y&=\Delta\sin\alpha\cos\ObservedBeta\\
z'-z&=\Delta\sin\ObservedBeta
\end{aligned}
\]
so dass $\Delta$ der Abstand des Planeten von der Erde, $\ObservedBeta$ die Neigung der von dem
Planeten zur Erde gezogenen geraden Linie gegen eine, parallel mit der Ebene
der $z$, durch die Erde gelegte Ebene; $\alpha$ der Winkel der Projection jener geraden
Linie auf diese Ebene gegen eine, parallel mit der Ebene der $y$, durch die
Erde gelegten Ebene sein werden. Auf der Himmelskugel bestimmen also $\alpha$ und
$\ObservedBeta$ die Lage des geocentrischen Orts des Planeten gegen die von den Ebenen der
$z$, $y$, $x$ (oder vielmehr ihnen parallel durch die Erde gelegten) gebildeten grössten
Kreise ganz eben so wie Länge und Breite die Lage gegen die Ecliptik und die
Coluren der Nachtgleichen und Sonnenwenden.

Vermöge der Beschaffenheit der Bahn des Planeten wird man zwischen
$x$, $y$, $z$ zwei Gleichungen haben, daher man diese drei veränderlichen Grössen
als Functionen \textit{einer} ansehen kann, die wir durch $t$ bezeichnen und übrigens noch
unbestimmt lassen wollen. Eben so sollen $x'$, $y'$, $z'$ Functionen der veränderli\SourceLeaf{113}%
chen Grösse $t'$ sein. Es sind also $\alpha$ und $\ObservedBeta$ Functionen von beiden $t$, $t'$, die
durch die Differentialgleichungen
\[
\mathrm{d}\alpha=p\mathrm{d}t+p'\mathrm{d}t',\qquad \mathrm{d}\ObservedBeta=q\mathrm{d}t+q'\mathrm{d}t'
\]
bestimmt werden mögen.

Dies vorausgesetzt, ist offenbar, dass, wenn man $t$, $t'$ sich zugleich so ändern
lässt, dass $\mathrm{d}t:\mathrm{d}t'=-p':p$, dadurch $\alpha$ ungeändert bleibe, $\ObservedBeta$ aber so
lange zu oder abnehmen werde, bis es einen grössten oder kleinsten Werth erreicht
hat. Dies geschieht offenbar, wenn $pq'-qp'=0$ wird. Nun ist klar,
dass die Combination \textit{aller} Werthe von $t$ und $t'$ alle mögliche geocentrische Örter
des Planeten gibt; und dass von allen solchen Combinationen, die einerlei $\alpha$
geben, diejenige Statt finden muss, wo $\ObservedBeta$ ein Grösstes oder Kleinstes wird, wenn
der geocentrische Ort in die Grenzen des Zodiacus des Planeten fallen soll. Hieraus
folgt also, dass diese Grenzen durch die Bedingungsgleichung $pq'-qp'=0$
bestimmt werden.

Die Differentiation obiger Gleichungen gibt
\[
\begin{aligned}
\mathrm{d}x'-\mathrm{d}x&=\cos\alpha\cos\ObservedBeta\,\mathrm{d}\Delta-\Delta\sin\alpha\cos\ObservedBeta\,\mathrm{d}\alpha-\Delta\cos\alpha\sin\ObservedBeta\,\mathrm{d}\ObservedBeta\\
\mathrm{d}y'-\mathrm{d}y&=\sin\alpha\cos\ObservedBeta\,\mathrm{d}\Delta+\Delta\cos\alpha\cos\ObservedBeta\,\mathrm{d}\alpha-\Delta\sin\alpha\sin\ObservedBeta\,\mathrm{d}\ObservedBeta\\
\mathrm{d}z'-\mathrm{d}z&=\sin\ObservedBeta\,\mathrm{d}\Delta+\Delta\cos\ObservedBeta\,\mathrm{d}\ObservedBeta
\end{aligned}
\]

Hieraus folgt leicht in Verbindung mit jenen Gleichungen
\[
\begin{gathered}
-\sin\alpha(\mathrm{d}x'-\mathrm{d}x)+\cos\alpha(\mathrm{d}y'-\mathrm{d}y)=\Delta\cos\ObservedBeta\,\mathrm{d}\alpha\\
\begin{aligned}
-\cos\alpha\sin\ObservedBeta(\mathrm{d}x'-\mathrm{d}x)-\sin\alpha\sin\ObservedBeta(\mathrm{d}y'-\mathrm{d}y)&\\
{}+\cos\ObservedBeta(\mathrm{d}z'-\mathrm{d}z)&=\Delta\mathrm{d}\ObservedBeta
\end{aligned}
\end{gathered}
\]

Es wird also, vermöge der partiellen Differentialien
\[
\begin{aligned}
\Delta\cos\ObservedBeta\,p\mathrm{d}t&=\sin\alpha\,\mathrm{d}x-\cos\alpha\,\mathrm{d}y\\
\Delta\cos\ObservedBeta\,p'\mathrm{d}t'&=-\sin\alpha\,\mathrm{d}x'+\cos\alpha\,\mathrm{d}y'\\
\Delta q\mathrm{d}t&=\cos\alpha\sin\ObservedBeta\,\mathrm{d}x+\sin\alpha\sin\ObservedBeta\,\mathrm{d}y-\cos\ObservedBeta\,\mathrm{d}z\\
\Delta q'\mathrm{d}t'&=-\cos\alpha\sin\ObservedBeta\,\mathrm{d}x'-\sin\alpha\sin\ObservedBeta\,\mathrm{d}y'+\cos\ObservedBeta\,\mathrm{d}z'
\end{aligned}
\]

Diese Werthe von $p$, $p'$, $q$, $q'$ in der Bedingungsgleichung $pq'=p'q$ substituirt,
wird nach den gehörigen Reductionen
\[
\begin{aligned}
\cos\alpha\cos\ObservedBeta(\mathrm{d}y'\mathrm{d}z-\mathrm{d}y\,\mathrm{d}z')&+\sin\alpha\cos\ObservedBeta(\mathrm{d}z'\mathrm{d}x-\mathrm{d}z\,\mathrm{d}x')\\
&+\sin\ObservedBeta(\mathrm{d}x'\mathrm{d}y-\mathrm{d}x\,\mathrm{d}y')=0
\end{aligned}
\]
\SourceLeaf{114}%
oder wenn man mit $\Delta$ multiplicirt
\[
\begin{aligned}
(x'-x)(\mathrm{d}y'\mathrm{d}z-\mathrm{d}y\,\mathrm{d}z')&+(y'-y)(\mathrm{d}z'\mathrm{d}x-\mathrm{d}z\,\mathrm{d}x')\\
&+(z'-z)(\mathrm{d}x'\mathrm{d}y-\mathrm{d}x\,\mathrm{d}y')=0
\end{aligned}
\]
welche Gleichung sich noch besser in folgender Form darstellen lässt:
\[
\left.\begin{aligned}
\mathrm{d}x(y'\mathrm{d}z'-z'\mathrm{d}y')&+\mathrm{d}x'(y\,\mathrm{d}z-z\,\mathrm{d}y)\\
{}+\mathrm{d}y(z'\mathrm{d}x'-x'\mathrm{d}z')&+\mathrm{d}y'(z\,\mathrm{d}x-x\,\mathrm{d}z)\\
{}+\mathrm{d}z(x'\mathrm{d}y'-y'\mathrm{d}x')&+\mathrm{d}z'(x\,\mathrm{d}y-y\,\mathrm{d}x)
\end{aligned}\right\}=0
\]

Diese Gleichung enthält allgemein die Relation zwischen den Örtern der
Erde und des Planeten, bei welchen der geocentrische Ort des letztern in die
Grenzen fällt, und man darf darin nur für $x$, $y$, $z$ ihre Werthe durch $t$, und
für $x'$, $y'$, $z'$ ihre Werthe durch $t'$, nach Beschaffenheit der Bahn, substituiren,
um eine endliche Gleichung zwischen $t$ und $t'$ zu erhalten. Es schien der Mühe
werth, jene Gleichung durch eine allgemeine Analyse zu entwickeln; übrigens
aber ist es nicht schwer zu zeigen, dass sie zugleich die Bedingungsgleichung sei,
\textit{dass die Tangenten an den Örtern der Erde und des Planeten in Einer Ebene liegen},
und gerade diese Bedingung aus den Erfordernissen unserer Aufgabe abzuleiten.
Kürze halber halten wir uns indessen hiebei nicht länger auf.

Für die Grössen $t$, $t'$, die wir bisher unbestimmt gelassen haben, nehmen
wir am bequemsten die heliocentrischen Winkel-Abstände des Planeten und der
Erde in ihren Bahnen von der gemeinschaftlichen Knotenlinie (und zwar vom
aufsteigenden Knoten der Planeten-Bahn auf der Erd-Bahn). Bezeichnen wir
nun die Entfernungen des Planeten und der Erde von der Sonne durch $r$, $r'$, so
werden sich die Coordinaten auf folgende Art ausdrücken lassen:
\[
\begin{aligned}
x&=r\sin a\sin(t+A)\\
y&=r\sin b\sin(t+B)\\
z&=r\sin c\sin(t+C)\\
x'&=r'\sin a'\sin(t'+A')\\
y'&=r'\sin b'\sin(t'+B')\\
z'&=r'\sin c'\sin(t'+C')
\end{aligned}
\]

Hievon, so wie von der Bedeutung der Constanten $a$, $A$ u.\,s.\,w. wird man
sich leicht durch Generalisirung der in der \textit{Mon. Corr.} Mai 1804 [S.\,98 d.\,B.]
vorgetragenen Untersuchung Rechenschaft geben können. Sind nun ferner $k$, $k'$
\SourceLeaf{115}%
die halben Parameter der Kegelschnitte. welche der Planet und die Erde beschreiben;
$e$, $e'$ die Excentricitäten; $g$, $g'$ die Winkel-Abstände der Sonnenfernen
von der Knotenlinie, so wird
\[
r=\frac{k}{1-e\cos(t-g)},\qquad r'=\frac{k'}{1-e'\cos(t'-g')}
\]

Hieraus findet sich nach gehöriger Rechnung
\[
\mathrm{d}x=\frac{k\sin a\,\mathrm{d}t}{(1-e\cos(t-g))^2}\times(\cos(t+A)-e\cos(g+A))
\]

Die Werthe von $\mathrm{d}y$, $\mathrm{d}z$ haben eine ähnliche Gestalt, und man braucht,
um sie zu erhalten, nur $a$, $A$ mit $b$, $B$ oder mit $c$, $C$ zu vertauschen. Die
Werthe von $\mathrm{d}x'$, $\mathrm{d}y'$, $\mathrm{d}z'$ erhält man aus denen von $\mathrm{d}x$, $\mathrm{d}y$, $\mathrm{d}z$, wenn man statt
der auf den Planeten sich beziehenden Grössen die analogen für die Erde setzt.

Die Entwickelung von $y\,\mathrm{d}z-z\,\mathrm{d}y$ geschieht bequemer aus den Werthen
von $y$, $z$, ehe man darin den Werth von $r$ substituirt hat: man erhält so
\[
y\,\mathrm{d}z-z\,\mathrm{d}y=rr\sin b\sin c\sin(B-C)\mathrm{d}t=rr\cos a\,\mathrm{d}t
\]
(man sehe den angeführten Aufsatz [S.\,100 d.\,B.]), und eben so
\[
\begin{aligned}
z\,\mathrm{d}x-x\,\mathrm{d}z&=rr\cos b\,\mathrm{d}t\\
x\,\mathrm{d}y-y\,\mathrm{d}x&=rr\cos c\,\mathrm{d}t
\end{aligned}
\]

Ganz ähnliche Werthe finden sich für die drei analogen, auf die Erde Bezug habenden
Ausdrücke.

Durch Substitution aller dieser Werthe wird die obige Bedingungsgleichung
nach den gehörigen Reductionen folgende
\[
\begin{aligned}
&k'\cos a'\sin a\,(\cos(t+A)-e\cos(g+A))\\
&{}+k'\cos b'\sin b\,(\cos(t+B)-e\cos(g+B))\\
&{}+k'\cos c'\sin c\,(\cos(t+C)-e\cos(g+C))\\
&{}+k\cos a\sin a'\,(\cos(t'+A')-e'\cos(g'+A'))\\
&{}+k\cos b\sin b'\,(\cos(t'+B')-e'\cos(g'+B'))\\
&{}+k\cos c\sin c'\,(\cos(t'+C')-e'\cos(g'+C'))=0
\end{aligned}
\]

Durch zweckmässige Reductionen lassen sich die drei ersten Theile dieser
Gleichung, wenn man die Neigung der Planetenbahn gegen die Erdbahn durch
$i$ bezeichnet, in $k'\sin i\,(\cos t-e\cos g)$, die drei letzten in $-k\sin i\,(\cos t'-e'\cos g')$
\SourceLeaf{116}%
verwandeln. Wir können indessen der Mühe, diese an sich zwar nicht schwierigen,
aber doch etwas weitläuftigen Reductionen zu entwickeln, hier um so eher
überhoben sein, da wir zu demselben Resultate viel bequemer gelangen können,
wenn wir die drei bisher unbestimmt gelassenen Fundamental-Ebenen, auf die
sich die Coordinaten beziehen, auf eine zweckmässige Art bestimmen. Wir wollen
nemlich für die Ebene der $z$ die Ecliptik, und die Ebenen der $x$, $y$ so annehmen,
dass der Pol der erstern in den aufsteigenden Knoten der Planeten-Bahn,
der Pol der zweiten hingegen $90^0$ weiter vorwärts in der Ecliptik falle.
Auf der Seite \textit{dieser} Pole, so wie auf der Nordseite der Ecliptik, sollen die Coordinaten
$x$, $y$, $z$ positiv gesetzt werden. Es ist leicht zu übersehen, dass unter
diesen Voraussetzungen
\[
\begin{array}{llll}
a=90^0,&b=90^0+i,&c=i,&C=0\\
a'=90^0,&b'=90^0.&c'=0,&B'=0
\end{array}
\]
werde, und mithin die obige Gleichung in folgende übergehe:
\[
k'\sin i\,(\cos t-e\cos g)-k\sin i\,(\cos t'-e'\cos g')=0
\]
oder
\[
k'(\cos t-e\cos g)=k(\cos t'-e'\cos g')
\]

Aus der Theorie der Kegelschnitte lässt sich leicht zeigen, dass $\frac{k}{\cos t-e\cos g}$
und $\frac{k'}{\cos t'-e'\cos g'}$ den Abstand zwischen der Sonne und den Durchschnittspunkten
der Knotenlinie der beiden Bahnen mit den Tangenten an den Örtern des Planeten
und der Erde ausdrücken. Die eben gefundene Gleichung zeigt daher an,
dass diese beiden Tangenten die Knotenlinie in einem und demselben Punkte
schneiden, welches mit der oben berührten Bedingung übereinkommt, nach der
sie in einer und derselben Ebene liegen sollen.

In Ansehung der Lage der Planetenbahn gegen die Erdbahn sind drei Fälle
zu unterscheiden. Entweder schliesst jene diese ein, oder diese jene, oder beide
einander (gleich Kettenringen). Der erste Fall findet Statt, wenn der Planet in
der Knotenlinie auf beiden Seiten weiter von der Sonne absteht, als die Erde;
der zweite, wenn diese auf beiden Seiten weiter absteht, als der Planet; der
dritte, wenn auf einer Seite der Planet, auf der andern die Erde weiter von der
Sonne entfernt ist. Von den bisher bekannten Planeten hat keiner eine solche
Lage gegen die Erde oder gegen einen andern Planeten, wie der dritte Fall er\SourceLeaf{117}%
fordert; Cometen der Art aber gibts in Menge. Die analytische Bedingung für
den ersten Fall ist, dass $k-k'$ positiv, und, ohne Rücksicht auf das Zeichen,
grösser sei, als $k'e\cos g-ke'\cos g'$; für den zweiten, dass $k'-k$ diese Eigenschaften
habe; für den dritten, dass $k-k'$ oder $k'-k$, ohne Rücksicht auf das
Zeichen, kleiner als $k'e\cos g-ke'\cos g'$ sei.

In dem ersten dieser drei Fälle erhält man aus obiger Gleichung für jeden
beliebigen Werth von $t$, von $0^0$ bis $360^0$, zwei Werthe von $t'$; im zweiten gibt
jeder Werth von $t'$ zwei von $t$. Der eine Werth von $t'$ (im ersten, oder von $t$
im zweiten Falle) wird nemlich allemal zwischen $0$ und $180$, der andere zwischen
$180$ und $360^0$ liegen, oder vielmehr innerhalb noch engerer Grenzen, deren
nähere Bestimmung keine Schwierigkeiten hat. Im ersten Falle also entsprechen
jedem heliocentrischen Orte des Planeten zwei heliocentrische Örter
der Erde, aber nicht umgekehrt, sondern nur in zwei von einander getrennten
Stücken der Erdbahn (wovon das eine unterhalb oder südlich von der Ebene der
Planetenbahn, das andere oberhalb oder nördlich von derselben liegt) kann die
Erde den Planeten in seinen Grenzen sehen, und zwar in der nördlichen Grenze
nur, wenn $t'$ einen von seinen möglichen Werthen zwischen $0$ und $180$, in der
südlichen, wenn es einen zwischen $180$ und $360^0$ erhält. Eben so entsprechen
im zweiten Falle jedem heliocentrischen Orte der Erde zwei des Planeten, aber
nicht umgekehrt: sondern nur in zwei von einander getrennten Stücken seiner
Bahn, wovon das eine nördlich, das andere südlich von der Ecliptik liegt, kann
er der Erde in seinen Grenzen erscheinen, nemlich an der nördlichen für die zwischen
$0$ und $180$, in der südlichen für die zwischen $180$ und $360^0$ liegenden
Werthe.

Hingegen können im dritten Falle weder $t$ noch $t'$ alle, sondern nur zwischen
gewissen Grenzen liegende Werthe erhalten; oder sowohl die Erde als der
Planet müssen jedes in einem bestimmten Stücke seiner Bahn sein, wenn obige
Bedingungsgleichung Statt haben soll.

Hieraus ergibt sich nun, dass die geocentrischen Örter des Planeten, die
aus allen möglichen, obiger Gleichung Genüge thuenden Combinationen zwischen
den heliocentrischen Örtern des Planeten und der Erde entspringen, im ersten
und zweiten Falle zwei von einander getrennte in sich selbst zurücklaufende Linien
auf der Himmelskugel bilden, zwischen denen im ersten Falle der die Ebene
der Planetenbahn vorstellende grösste Kreis, im zweiten die Ecliptik liegt; im
\SourceLeaf{118}%
dritten Falle hingegen bilden jene geocentrischen Örter (wie die nähere Betrachtung
des Falles ohne Mühe zeigt) nur \textit{eine} in sich zurückkehrende Linie.

Den vorhergehenden Untersuchungen zu Folge kann nun der Zodiacus des
Planeten keine andere Grenzen haben als eben diese Linien. Es scheint daher
natürlich, zu schliessen, dass in den beiden ersten Fällen der Zodiacus des Planeten
die zwischen jenen beiden Linien liegende Zone, und im dritten einer von
von den beiden Räumen sei, in welche jene Linie die ganze Kugelfläche scheidet.
Allein dieser Schluss würde für die beiden ersten Fälle nicht immer, und
für den dritten nie richtig sein. Man darf nemlich hier (so wie in vielen andern
Fällen beim Gebrauch der Analyse, wo man eine ähnliche Vorsicht nicht immer
genug beobachtet) nicht vergessen, dass unsere Schlussfolge sich ganz auf die
Voraussetzung gründet, dass der Zodiacus des Planeten wirklich beschränkt sei,
und dass dieser von der Erde aus nicht in jedem Punkte des Himmels erscheinen
könne. Diese Voraussetzung findet aber, wie sich schon aus Gründen der \textit{Geometrie
der Lage} darthun lässt, in dem dritten Falle nicht Statt, und die gefundene
Linie kann also hier nicht die Grenze des Planeten-Zodiacus sein, da dieser
den \textit{ganzen} Himmel einnimmt. Im ersten und zweiten Falle aber wird es
zwar allemal wenigstens auf einer Seite der gefundenen Zone Stellen am Himmel
geben, wo der Planet nie erscheinen kann, und folglich gewiss die eine Linie
eine Grenze sein; allein demungeachtet kann es sich ereignen, dass es \textit{nur} auf
einer Seite solche ausgeschlossene Stellen gibt, daher dann die andere Linie keine
Grenze abgibt, sondern der dadurch von der Zone abgeschiedene Raum des Himmels
eben so gut ganz zum Zodiacus gehört, als die Zone selbst. Indessen
ist es hier nicht der Ort, diese Untersuchung vollständig auszuführen, und es
zu entwickeln, was denn in solchen Fällen jene Linien, da sie keine Grenzen
sind, eigentlich bedeuten. Hier können wir uns um so eher begnügen, die
Freunde der Analyse auf diese paradox scheinenden Phänomene aufmerksam gemacht
zu haben, da es sich leicht zeigen lässt, dass alle bis jetzt bekannten Planeten,
die hier zunächst unser Augenmerk sind, nie weder im Nordpol noch im
Südpol der Ecliptik von der Erde aus erscheinen, und folglich die erwähnten
Ausnahmen dabei nicht Statt haben können; daher ihr Zodiacus wirkliche \textit{Zonen},
und die beiden gefundenen Linien ihre Grenzen sein müssen.

Wir wollen nun noch theils zur weitern Erläuterung, theils des practischen
Gebrauchs wegen unsere Resultate auf die \textit{Pallas} und \textit{Ceres} anwenden, und die
\SourceLeaf{119}%
Grenzen ihrer Zodiacus so abstecken, dass man sie danach in die Sternkarten eintragen
könne. Wegen der Perturbationen werden zwar diese Grenzen noch einiger
Erweiterung, und wegen der Veränderung, die die Elemente in Zukunft
noch erleiden werden, einiger Änderungen bedürfen; allein in practischer Rücksicht
werden dieselben unerheblich, und, wenn die Beobachtungen nach Jahren
sie merklich machen werden, eben darum sogar interessant sein, weil sie dann
die nach und nach eintretende Unzulänglichkeit der hier zum Grunde gelegten
elliptischen Elemente auf eine in die Augen fallende Art zeigen werden.

Für die \textit{Pallas} setzen wir nach den neuesten Elementen für den Anfang
von 1803 (\textit{Mon. Corr.} 1804 März.)
\[
\begin{array}{rl@{\qquad}rl}
e&=0.2457396&\varphi&=14^0\;13'\;31''966\\
k&=2.602122&&\\
g&=128^0\;49'\;20''7&&\\
\AscendingNode&=172\;28\;13.7&i&=34\;38\;1.1
\end{array}
\]

Für die Erde hingegen
\[
\begin{array}{rl@{\qquad}rl}
e'&=0.016792&\varphi&=0^0\;57'\;43''802\\
k'&=0.999718&&\\
g'&=107^0\;5'\;39''&&\\
\mathrm{Aph.}&=279\;33\;52.7&&
\end{array}
\]

Nach Substitution dieser Werthe wird unsere obige Gleichung
\[
\cos t'=0.384193\cos t+0.0542514
\]

Hieraus folgt, dass die beiden äussersten Werthe von $\cos t'$ diese sind
$+0.438444$ und $-0.329942$; es liegen also alle möglichen Werthe von $t'$ einerseits
zwischen $63^0\;59'\;43''$ und $109^0\;15'\;55''$; andererseits zwischen $250^0\;44'\;5''$
und $296^0\;0'\;17''$; daher die \textit{Pallas} der Erde nur dann in ihren Grenzen erscheinen
kann, wenn die heliocentrische Länge jener zwischen $236^0\;28'$ und $281^0\;44'$
oder zwischen $63^0\;12'$ und $180^0\;29'$ fällt; also etwa vom 18.\,Mai bis 4.\,Juli, und
vom 26.\,November bis 9.\,Januar. In dem ersten Theile ihrer Bahn befindet sich
die Erde südlich, im andern nördlich von der Ebene der Pallas-Bahn; daher ihr
in jenem die \textit{Pallas} an der nördlichen, in diesem an der südlichen Grenze ihres
Zodiacus erscheinen wird. Auch ist es nicht schwer, zu zeigen, dass die \textit{Pallas}
\SourceLeaf{120}%
jedes Jahr zu den bestimmten Zeiten einmal die nördliche und einmal die südliche
Grenze streifen muss. — Um nun eine hinlängliche Anzahl von Punkten
aus beiden Grenzen zu erhalten, wollen wir für $t$ der Reihe nach alle Werthe
von $0$ bis $360^0$ von $10$ zu $10$ Grad annehmen, und aus der Verbindung jedes
derselben mit den beiden zugehörigen, aus obiger Formel zu bestimmenden Werthen
von $t'$ die entsprechenden geocentrischen Örter sogleich in Rectascension
und Declination ableiten, in welcher Absicht das in dem oben erwähnten Aufsatze
erklärte Verfahren und die dabei bereits berechneten Constanten angewandt
werden können. Die Resultate dieser Rechnungen stellt folgende Tafel dar:

\begin{center}
\begingroup\small\setlength{\tabcolsep}{5pt}
\begin{tabular}{r||rr|rrl||rr|rrl}
&\multicolumn{5}{c||}{Nördliche Grenze}&\multicolumn{5}{c}{Südliche Grenze}\\
$t$&\multicolumn{2}{c|}{Gerade Aufst.}&\multicolumn{3}{c||}{Abweichung}&\multicolumn{2}{c|}{Gerade Aufst.}&\multicolumn{3}{c}{Abweichung.}\\\hline
$0^0$&$148^0$&$4'$&$12^0$&$56'$&nördl.&$196^0$&$54'$&$7^0$&$12'$&südl.\\
$10$&$159$&$30$&$15$&$56$&&$205$&$54$&$4$&$26$&\\
$20$&$171$&$54$&$18$&$55$&&$214$&$15$&$2$&$11$&\\
$30$&$185$&$6$&$21$&$38$&&$222$&$3$&$0$&$23$&\\
$40$&$198$&$48$&$23$&$52$&&$229$&$27$&$1$&$1$&nördl.\\
$50$&$212$&$41$&$25$&$26$&&$236$&$32$&$2$&$4$&\\
$60$&$226$&$28$&$26$&$16$&&$243$&$26$&$2$&$48$&\\
$70$&$239$&$59$&$26$&$23$&&$250$&$16$&$3$&$16$&\\
$80$&$253$&$8$&$25$&$53$&&$257$&$6$&$3$&$28$&\\
$90$&$265$&$54$&$24$&$49$&&$264$&$1$&$3$&$25$&\\
$100$&$278$&$18$&$23$&$19$&&$271$&$4$&$3$&$5$&\\
$110$&$290$&$25$&$21$&$27$&&$278$&$18$&$2$&$30$&\\
$120$&$302$&$19$&$19$&$17$&&$285$&$44$&$1$&$37$&\\
$130$&$314$&$4$&$16$&$53$&&$293$&$23$&$0$&$27$&\\
$140$&$325$&$43$&$14$&$22$&&$301$&$16$&$1$&$2$&südl.\\
$150$&$337$&$16$&$11$&$47$&&$309$&$24$&$2$&$50$&\\
$160$&$348$&$42$&$9$&$14$&&$317$&$49$&$4$&$57$&\\
$170$&$359$&$55$&$6$&$49$&&$326$&$34$&$7$&$23$&\\
$180$&$10$&$49$&$4$&$40$&&$335$&$41$&$10$&$8$&\\
$190$&$21$&$15$&$2$&$49$&&$345$&$12$&$13$&$10$&\\
$200$&$31$&$5$&$1$&$21$&&$355$&$12$&$16$&$24$&\\
$210$&$40$&$14$&$0$&$14$&&$5$&$44$&$19$&$53$&\\
$220$&$48$&$39$&$0$&$31$&südl.&$16$&$53$&$23$&$6$&\\
$230$&$56$&$23$&$0$&$59$&&$28$&$41$&$26$&$16$&\\
$240$&$63$&$28$&$1$&$10$&&$41$&$12$&$29$&$5$&\\
$250$&$70$&$2$&$1$&$9$&&$54$&$25$&$31$&$21$&\\
$260$&$76$&$12$&$0$&$57$&&$68$&$16$&$32$&$52$&\\
$270$&$82$&$5$&$0$&$35$&&$82$&$36$&$33$&$29$&\\
$280$&$87$&$51$&$0$&$4$&&$97$&$12$&$33$&$5$&\\
$290$&$93$&$38$&$0$&$37$&nördl.&$111$&$47$&$31$&$39$&\\
$300$&$99$&$35$&$1$&$30$&&$126$&$6$&$29$&$14$&\\
$310$&$105$&$53$&$2$&$36$&&$139$&$52$&$26$&$0$&\\
$320$&$112$&$40$&$3$&$59$&&$152$&$57$&$22$&$12$&\\
$330$&$120$&$8$&$5$&$40$&&$165$&$12$&$18$&$9$&\\
$340$&$128$&$26$&$7$&$44$&&$176$&$37$&$14$&$10$&\\
$350$&$137$&$43$&$10$&$10$&&$187$&$10$&$10$&$28$&\\
$360$&$148$&$4$&$12$&$56$&&$196$&$54$&$7$&$12$&
\end{tabular}
\endgroup
\end{center}
\SourceLeaf{121}%
Zu grösserer Bequemlichkeit sind durch schickliche Interpolations-Methoden
zwischen diese $72$ Punkte folgende $144$ eingeschaltet, bei denen die Rectascensionen
von $5$ zu $5$ Grad zunehmen.

\begin{center}
\begingroup\small\setlength{\tabcolsep}{3pt}
\begin{tabular}{r|rrl|rrl||r|rrl|rrl}
\multicolumn{14}{c}{\textit{Zodiacus der Pallas}}\\[4pt]
&\multicolumn{6}{c||}{Abweichung}&&\multicolumn{6}{c}{Abweichung}\\
\shortstack{Gerade\\Aufst.}&\multicolumn{3}{c|}{\shortstack{der nördlichen\\Grenze}}&\multicolumn{3}{c||}{\shortstack{der südlichen\\Grenze}}&\shortstack{Gerade\\Aufst.}&\multicolumn{3}{c|}{\shortstack{der nördlichen\\Grenze}}&\multicolumn{3}{c}{\shortstack{der südlichen\\Grenze}}\\\hline
$0^0$&$6^0$&$48'$&nördl.&$17^0$&$56'$&südl.&$180^0$&$20^0$&$38'$&nördl.&$12^0$&$58'$&südl.\\
$5$&$5$&$47$&&$19$&$32$&&$185$&$21$&$37$&&$11$&$12$&\\
$10$&$4$&$49$&&$21$&$4$&&$190$&$22$&$31$&&$9$&$29$&\\
$15$&$3$&$54$&&$22$&$33$&&$195$&$23$&$19$&&$7$&$49$&\\
$20$&$3$&$2$&&$23$&$59$&&$200$&$24$&$2$&&$6$&$13$&\\
$25$&$2$&$13$&&$25$&$20$&&$205$&$24$&$39$&&$4$&$42$&\\
$30$&$1$&$29$&&$26$&$36$&&$210$&$25$&$11$&&$3$&$17$&\\
$35$&$0$&$50$&&$27$&$46$&&$215$&$25$&$37$&&$2$&$0$&\\
$40$&$0$&$16$&&$28$&$51$&&$220$&$25$&$58$&&$0$&$49$&\\
$45$&$0$&$14$&südl.&$29$&$49$&&$225$&$26$&$12$&&$0$&$13$&nördl.\\
$50$&$0$&$37$&&$30$&$41$&&$230$&$26$&$22$&&$1$&$7$&\\
$55$&$0$&$55$&&$31$&$26$&&$235$&$26$&$25$&&$1$&$52$&\\
$60$&$1$&$6$&&$32$&$4$&&$240$&$26$&$23$&&$2$&$28$&\\
$65$&$1$&$11$&&$32$&$36$&&$245$&$26$&$16$&&$2$&$56$&\\
$70$&$1$&$9$&&$33$&$0$&&$250$&$26$&$3$&&$3$&$15$&\\
$75$&$1$&$0$&&$33$&$17$&&$255$&$25$&$45$&&$3$&$26$&\\
$80$&$0$&$44$&&$33$&$27$&&$260$&$25$&$23$&&$3$&$28$&\\
$85$&$0$&$21$&&$33$&$30$&&$265$&$24$&$55$&&$3$&$23$&\\
$90$&$0$&$10$&nördl.&$33$&$25$&&$270$&$24$&$22$&&$3$&$9$&\\
$95$&$0$&$49$&&$33$&$13$&&$275$&$23$&$46$&&$2$&$48$&\\
$100$&$1$&$34$&&$32$&$54$&&$280$&$23$&$5$&&$2$&$19$&\\
$105$&$2$&$27$&&$32$&$27$&&$285$&$22$&$20$&&$1$&$43$&\\
$110$&$3$&$25$&&$31$&$53$&&$290$&$21$&$31$&&$1$&$0$&\\
$115$&$4$&$29$&&$31$&$11$&&$295$&$20$&$39$&&$0$&$10$&\\
$120$&$5$&$38$&&$30$&$23$&&$300$&$19$&$43$&&$0$&$46$&südl.\\
$125$&$6$&$52$&&$29$&$27$&&$305$&$18$&$45$&&$1$&$49$&\\
$130$&$8$&$9$&&$28$&$24$&&$310$&$17$&$44$&&$2$&$58$&\\
$135$&$9$&$27$&&$27$&$14$&&$315$&$16$&$42$&&$4$&$12$&\\
$140$&$10$&$47$&&$25$&$57$&&$320$&$15$&$37$&&$5$&$32$&\\
$145$&$12$&$7$&&$24$&$35$&&$325$&$14$&$31$&&$6$&$56$&\\
$150$&$13$&$27$&&$23$&$7$&&$330$&$13$&$24$&&$8$&$24$&\\
$155$&$14$&$46$&&$21$&$33$&&$335$&$12$&$17$&&$9$&$56$&\\
$160$&$16$&$4$&&$19$&$55$&&$340$&$11$&$10$&&$11$&$30$&\\
$165$&$17$&$19$&&$18$&$14$&&$345$&$10$&$3$&&$13$&$6$&\\
$170$&$18$&$29$&&$16$&$30$&&$350$&$8$&$57$&&$14$&$43$&\\
$175$&$19$&$36$&&$14$&$44$&&$355$&$7$&$52$&&$16$&$20$&\\
$180$&$20$&$38$&&$12$&$58$&&$360$&$6$&$48$&&$17$&$56$&
\end{tabular}
\endgroup
\end{center}

Für die \textit{Ceres} haben wir, nach den letzten Elementen für 1803
\[
\begin{aligned}
e&=0.0788941.\\
k&=2.750681\\
g&=245^0\;34'\;56''
\end{aligned}
\]
\SourceLeaf{122}
für die Erde $e'$ und $k'$ wie oben, $g'=198^0\;35'\;31''$.
Hieraus wird die Bedingungsgleichung
\[
\cos t'=.\,0.363444\cos t-0.004063
\]

Die Werthe von $t'$ liegen also von $68^0\;56'\;16''$ bis $111^0\;33'\;43''$ und von
$248^0\;26'\;17''$ bis $291^0\;3'\;44''$, folglich die heliocentrische Länge der Erde von
$149^0\;55'$ bis $192^0\;32'$ und von $329^0\;25'$ bis $12^0\;2'$; daher die Erde etwa nur
vom 19.\,Februar bis 3.\,April die \textit{Ceres} in ihren nördlichen, und vom 23.\,August
bis 6.\,October in ihren südlichen Grenzen sehen kann. Vermittelst dieser Formel sind,
eben so wie vorhin bei der \textit{Pallas}, 36 Punkte in jeder Grenze des Zodiacus der
\textit{Ceres} in Rectascension und Declination berechnet worden, wobei in Beziehung
auf den mehr erwähnten Aufsatz für die Constanten $a$, $A$ u.\,s.\,w. folgende Werthe
gebraucht sind:
\[
\begin{array}{rcrrr}
a&=&79^0&30'&5''\\
b&=&114&42&11\\
c&=&27&7&24\\
A&=&170&49&4\\
B&=&85&42&29\\
C&=&59&36&33
\end{array}
\]
so dass, wenn $v$ die wahre Anomalie der \textit{Ceres} bedeutet, die drei Coordinaten
durch folgende Formeln dargestellt werden:
\[
\begin{aligned}
x&=\frac{\alpha\sin(v+56^0\;24'\;0'')}{1-e\cos v}\\[8pt]
y&=\frac{\ObservedBeta\sin(v+331^0\;17'\;25'')}{1-e\cos v}\\[8pt]
z&=\frac{\gamma\sin(v+305^0\;11'\;29'')}{1-e\cos v}
\end{aligned}
\]
wo
\[
\begin{aligned}
\log\alpha&=\log k\sin a=0.432108\\
\log\ObservedBeta&=\log k\sin b=0.397758\\
\log\gamma&=\log k\sin c=0.098319
\end{aligned}
\]
\SourceLeaf{123}

Anstatt dieser $72$ Punkte begnügen wir uns hier damit, nur die auf ähnliche Art wie
bei der \textit{Pallas} zwischen dieselben eingeschalteten $144$ in folgender Tafel beizufügen:

\begin{center}
\begingroup\small\setlength{\tabcolsep}{3pt}
\begin{tabular}{r|rrl|rrl||r|rrl|rrl}
\multicolumn{14}{c}{\textit{Zodiacus der Ceres}}\\[4pt]
&\multicolumn{6}{c||}{Abweichung}&&\multicolumn{6}{c}{Abweichung}\\
\shortstack{Gerade\\Aufst.}&\multicolumn{3}{c|}{\shortstack{der nördlichen\\Grenze}}&\multicolumn{3}{c||}{\shortstack{der südlichen\\Grenze}}&\shortstack{Gerade\\Aufst.}&\multicolumn{3}{c|}{\shortstack{der nördlichen\\Grenze}}&\multicolumn{3}{c}{\shortstack{der südlichen\\Grenze}}\\\hline
$0^0$&$8^0$&$37'$&südl.&$17^0$&$29'$&südl.&$180^0$&$18^0$&$49'$&nördl.&$8^0$&$17'$&nördl.\\
$5$&$6$&$15$&&$15$&$13$&&$185$&$16$&$31$&&$5$&$56$&\\
$10$&$3$&$45$&&$12$&$51$&&$190$&$14$&$5$&&$3$&$28$&\\
$15$&$1$&$11$&&$10$&$26$&&$195$&$11$&$33$&&$0$&$55$&\\
$20$&$1$&$27$&nördl.&$7$&$56$&&$200$&$8$&$57$&&$1$&$41$&südl.\\
$25$&$4$&$7$&&$5$&$18$&&$205$&$6$&$16$&&$4$&$19$&\\
$30$&$6$&$48$&&$2$&$38$&&$210$&$3$&$34$&&$6$&$56$&\\
$35$&$9$&$26$&&$0$&$3$&&$215$&$0$&$52$&&$9$&$32$&\\
$40$&$12$&$1$&&$2$&$30$&nördl.&$220$&$1$&$47$&südl.&$12$&$3$&\\
$45$&$14$&$31$&&$4$&$58$&&$225$&$4$&$21$&&$14$&$29$&\\
$50$&$16$&$54$&&$7$&$20$&&$230$&$6$&$50$&&$16$&$48$&\\
$55$&$19$&$9$&&$9$&$35$&&$235$&$9$&$12$&&$18$&$59$&\\
$60$&$21$&$16$&&$11$&$42$&&$240$&$11$&$25$&&$21$&$1$&\\
$65$&$23$&$12$&&$13$&$39$&&$245$&$13$&$28$&&$22$&$53$&\\
$70$&$24$&$58$&&$15$&$27$&&$250$&$15$&$21$&&$24$&$35$&\\
$75$&$26$&$34$&&$17$&$4$&&$255$&$17$&$3$&&$26$&$6$&\\
$80$&$27$&$59$&&$18$&$30$&&$260$&$18$&$34$&&$27$&$26$&\\
$85$&$29$&$13$&&$19$&$45$&&$265$&$19$&$53$&&$28$&$36$&\\
$90$&$30$&$16$&&$20$&$49$&&$270$&$21$&$1$&&$29$&$35$&\\
$95$&$31$&$8$&&$21$&$42$&&$275$&$21$&$56$&&$30$&$24$&\\
$100$&$31$&$50$&&$22$&$24$&&$280$&$22$&$40$&&$31$&$1$&\\
$105$&$32$&$21$&&$22$&$55$&&$285$&$23$&$13$&&$31$&$29$&\\
$110$&$32$&$42$&&$23$&$14$&&$290$&$23$&$34$&&$31$&$46$&\\
$115$&$32$&$51$&&$23$&$22$&&$295$&$23$&$43$&&$31$&$52$&\\
$120$&$32$&$50$&&$23$&$19$&&$300$&$23$&$41$&&$31$&$48$&\\
$125$&$32$&$40$&&$23$&$5$&&$305$&$23$&$28$&&$31$&$34$&\\
$130$&$32$&$18$&&$22$&$40$&&$310$&$23$&$4$&&$31$&$9$&\\
$135$&$31$&$46$&&$22$&$3$&&$315$&$22$&$27$&&$30$&$34$&\\
$140$&$31$&$3$&&$21$&$16$&&$320$&$21$&$40$&&$29$&$49$&\\
$145$&$30$&$9$&&$20$&$17$&&$325$&$20$&$40$&&$28$&$53$&\\
$150$&$29$&$4$&&$19$&$6$&&$330$&$19$&$30$&&$27$&$46$&\\
$155$&$27$&$49$&&$17$&$45$&&$335$&$18$&$8$&&$26$&$29$&\\
$160$&$26$&$22$&&$16$&$12$&&$340$&$16$&$35$&&$25$&$1$&\\
$165$&$24$&$44$&&$14$&$28$&&$345$&$14$&$51$&&$23$&$23$&\\
$170$&$22$&$56$&&$12$&$34$&&$350$&$12$&$55$&&$21$&$34$&\\
$175$&$20$&$57$&&$10$&$30$&&$355$&$10$&$51$&&$19$&$36$&\\
$180$&$18$&$49$&&$8$&$17$&&$360$&$8$&$37$&&$17$&$29$&\\
\end{tabular}
\endgroup
\end{center}
